\setcounter{chapter}{1}\setcounter{section}{0}\setcounter{subsection}{0}
\chapter{Statistiques}
\label{statistiques}\hypertarget{statistiques}{}
\compteursfinauxexamens

\section{Statistiques générales des sujets}

\begin{footnotesize}{\textit{Sauf mention du contraire, on appelle dans le texte qui suit «~sujet~»{} l'ensemble composé d'un texte et de deux questions (l'une littéraire, l'autre philosophique). On appelle «~question~»{} les questions d'interprétation et essais qui sont proposés dans chacun des sujets, ce qui fait donc deux questions par sujet. On appelle «~Essai~»{}, conformément à la terminologie des épreuves de Terminale ce qui en classe de Première est appelé «~Question de réflexion~»{}.}}\end{footnotesize}

\vspace{3pt}

À ce jour (\today), \textbf{\arabic{nosujet}}~sujets proposés ont été recensés (dont 4 «~Sujets 0~»). Il y a eu \textbf{\arabic{sujetRS}}~sujets sur «~\RS{}~» et \textbf{\arabic{sujetHQ}} sur «~\HQ{}~». Parmi eux, il y a eu \arabic{nosujetintlitteraire}~questions d'interprétation littéraire et autants d'essais philosophiques, et \arabic{nosujetintphilosophique}~questions d'interprétation philosophique et autant d'essais littéraires.

Pour \textbf{«~\RS{}~»{}}, il y a eu :\setlength{\itemsep}{2pt}

\vspace*{-14pt}

\begin{itemize}
\item \arabic{sujetETA}~questions qui portaient exclusivement sur «\ETA{}»{} et \arabic{sujetallETA}~questions qui faisaient appel au chapitre sur « \ETA{} »{},
\item \arabic{sujetEXS}~questions qui portaient exclusivement sur « \EXS{}»{} et \arabic{sujetallEXS}~questions qui faisaient appel au chapitre sur « \EXS{}»{},
\item \arabic{sujetMM}~questions qui portaient exclusivement sur «\MM{}»{} et \arabic{sujetallMM}~questions qui faisaient appel au chapitre sur « \MM{} »{}. 
\end{itemize}

\vspace*{-14pt}

Parmi les questions qui font appel à plusieurs chapitres du semestre, il y a eu :

\vspace*{-14pt}

\begin{itemize}
\item \arabic{sujetEXSTA}~question qui portait sur «\EXSTA{}»{},
\item \arabic{sujetEXSMM}~questions qui portaient sur «\EXSMM{}»{}, 
\item \arabic{sujetETAMM}~questions qui portaient sur «{}\ETAMM{}»{}, \item et \arabic{sujetOdS}~questions qui étaient des «\OdS{}»{}.
\end{itemize}

\vspace*{-14pt}

Pour \textbf{«\HQ{}»{}}, il y a eu :

\vspace*{-14pt}

\begin{itemize}
\item \arabic{sujetCC}~questions qui portaient exclusivement sur «\CC{}»{} et \arabic{sujetallCC}~questions qui faisaient appel au chapitre sur « \CC{} »{},
\item \arabic{sujetHV}~questions qui portaient exclusivement sur «{}\HV{}»{} et \arabic{sujetallHV}~questions qui faisaient appel au chapitre sur « \HV{} »{},
\item \arabic{sujetHL}~questions qui portaient exclusivement sur «{}\HL{}»{} et \arabic{sujetallHL}~questions qui faisaient appel au chapitre sur « \HL{} »{}.
\end{itemize}

\vspace*{-14pt}

Parmi les~questions qui font appel à plusieurs chapitres du semestre, il y a eu :

\vspace*{-14pt}

\begin{itemize}
\item \arabic{sujetHVL}~questions qui portaient sur «\HVL{}»{},
\item \arabic{sujetHVCC}~questions qui portaient sur «\HVCC{}»{},
\item \arabic{sujetHLCC}~sujets qui portaient sur «\HLCC{}»{},
\item et \arabic{sujetOdH}~questions qui étaient des «{}\OdH{}»{}.
\end{itemize}

\vspace*{-14pt}

Les questions qui font appel à des chapitres appartenant à deux semestres différents (c'est-à-dire à «\RS»{} et à «\HQ{}»{}) se répartissent ainsi~: \arabic{sujetEXSHV}~questions sur «\EXSHV{}»{}, \arabic{sujetEXSHL}~questions sur «\EXSHL{}»{} et \arabic{sujetMMHV}~question sur «\MMHV»{}, \arabic{sujetETACC}~questions qui portaient sur « \ETACC{}»{}.

Enfin, certaines questions faisaient appel à des connaissances de la classe de Première, à savoir \arabic{sujetHVDMRC}~questions sur «\HVDMRC{}»{}, \arabic{sujetEXSPdP}~question(s) sur «\EXSPdP{}»{}, et \arabic{sujetHLHA} sur « \HLHA »{}. On obtient le tableau qui suit, qui rend probablement ces données plus lisibles.

\vspace*{-4pt}\renewcommand{\arraystretch}{1.3}
\setlength\extrarowheight{1pt}

\begin{table}[h]
\centering
\begin{tabularx}{\linewidth+6.8pt}{|M{3.5cm}M{3.5cm}M{0.6cm}|M{0.6cm}M{3.5cm}M{3.5cm}|}
\hline
\rowcolor[HTML]{C0C0C0}
\multicolumn{2}{|c}{\small \textbf{Chapitre(s) de la Recherche de soi}} & \multicolumn{1}{c|}{\makecell{{\footnotesize Nombre}\\{\footnotesize de quest.}}}&\multicolumn{1}{c}{\makecell{{\footnotesize Nombre}\\{\footnotesize de quest.}}}& \multicolumn{2}{c|}{\small  \textbf{Chapitre(s) de l'Humanité en question}}\\\hline
\multicolumn{2}{|c}{\makecell[cc]{\small{Éducation, transmission et émancipation}}}&\multicolumn{1}{c|}{\arabic{sujetallETA}} & \multicolumn{1}{c}{\arabic{sujetallCC}} & \multicolumn{2}{c|}{\makecell[cc]{Création, continuités et ruptures}}\\
\multicolumn{2}{|c}{\makecell[cc]{\textit{exclusivement}}}&\multicolumn{1}{c|}{\arabic{sujetETA}} & \multicolumn{1}{c}{\arabic{sujetCC}} & \multicolumn{2}{c|}{\makecell[cc]{\textit{exclusivement}}}\\
\hline
\rowcolor[HTML]{efefef}
\multicolumn{2}{|c}{Les expressions de la sensibilité} & \multicolumn{1}{c|}{\arabic{sujetallEXS}}&\multicolumn{1}{c}{\arabic{sujetallHV}} & \multicolumn{2}{c|}{Histoire et violence}\\
\rowcolor[HTML]{efefef}
\multicolumn{2}{|c}{\textit{exclusivement}}&\multicolumn{1}{c|}{\arabic{sujetEXS}} & \multicolumn{1}{c}{\arabic{sujetHV}} & \multicolumn{2}{c|}{\textit{exclusivement}}\\\hline
\multicolumn{2}{|c}{Les métamorphoses du moi} & \multicolumn{1}{c|}{\arabic{sujetallMM}} & \multicolumn{1}{c}{\arabic{sujetallHL}}&\multicolumn{2}{c|}{L'humain et ses limites}\\
\multicolumn{2}{|c}{\makecell[cc]{\textit{exclusivement}}}&\multicolumn{1}{c|}{\arabic{sujetMM}} & \multicolumn{1}{c}{\arabic{sujetHL}} & \multicolumn{2}{c|}{\makecell[cc]{\textit{exclusivement}}}\\\hline
\rowcolor[HTML]{C0C0C0}
\multicolumn{6}{|c|}{\cellcolor[HTML]{C0C0C0}\textbf{Questions sur plusieurs chapitres du même semestre}} \\ \hline
\multicolumn{2}{|c}{\makecell[cc]{Les expressions de la sensibilité et \\éducation, transmission et émancipation}}&\multicolumn{1}{c|}{\arabic{sujetEXSTA}} & \multicolumn{1}{c}{\arabic{sujetHVCC}}&\multicolumn{2}{c|}{\makecell[cc]{Histoire et violence et création,\\ continuités et ruptures}}\\ \hline
\multicolumn{2}{|c}{\makecell{Les expressions de la sensibilité\\et les métamorphoses du moi}}&\multicolumn{1}{c|}{\arabic{sujetEXSMM}} & \multicolumn{1}{c}{\arabic{sujetHVL}}&\multicolumn{2}{c|}{\makecell{Histoire et violence et\\l'humain et ses limites}}\\ \hline
\multicolumn{2}{|c}{\makecell{Éducation, transmission et émancipation\\et les métamorphoses du moi}} & \multicolumn{1}{c|}{\arabic{sujetETAMM}} & \multicolumn{1}{c}{\arabic{sujetHLCC}} & \multicolumn{2}{c|}{\makecell{L'humain et ses limites et création,\\ continuités et ruptures}} \\ \hline
\multicolumn{2}{|c}{\makecell{Sujets originaux sur\\la Recherche de soi}}&\multicolumn{1}{c|}{\arabic{sujetOdS}} & \multicolumn{1}{c}{\arabic{sujetOdH}}&\multicolumn{2}{c|}{\makecell{Sujets originaux sur\\l'Humanité en question}}\\ \hline
\multicolumn{1}{|c|}{\cellcolor[HTML]{C0C0C0}}&\multicolumn{1}{c}{\makecell[cc]{{\footnotesize Les expressions de la}\\{\footnotesize sensibilité et histoire}\\{\footnotesize et violence}}}&\multicolumn{1}{c|}{\arabic{sujetEXSHV}}&\multicolumn{1}{c}{\arabic{sujetEXSPdP}}&\multicolumn{1}{c}{\makecell[cc]{{\footnotesize Les expressions de la}\\{\footnotesize sensibilité et les}\\{\footnotesize pouvoirs de la parole}}}&\multicolumn{1}{|c|}{\cellcolor[HTML]{C0C0C0}}\\
\cline{2-5}\multicolumn{1}{|c|}{\cellcolor[HTML]{C0C0C0}}&\multicolumn{1}{c}{\makecell[cc]{{\footnotesize Les expressions de la}\\{\footnotesize sensibilité et l'humain}\\{\footnotesize et ses limites}}}&\multicolumn{1}{c|}{\arabic{sujetEXSHL}}&\multicolumn{1}{c}{\arabic{sujetHVDMRC}}&\multicolumn{1}{c}{\makecell[cc]{{\footnotesize \HV{} et}\\{\footnotesize découverte du monde \&}\\{\footnotesize rencontre des cultures}}}&\multicolumn{1}{|c|}{\cellcolor[HTML]{C0C0C0}}\\
\cline{2-5}\multicolumn{1}{|c|}{\cellcolor[HTML]{C0C0C0}}&\multicolumn{1}{c}{\makecell[cc]{{\footnotesize Éducation, transmission}\\{\footnotesize et émancipation et Cré-}\\{\footnotesize{-ation, continuités, etc.}}}}&\multicolumn{1}{c|}{\arabic{sujetETACC}}&\multicolumn{1}{c}{\arabic{sujetHLHA}}&\multicolumn{1}{c}{\makecell[cc]{{\footnotesize L'humain et ses limites}\\{\footnotesize et l'Homme et l'animal}}}&\multicolumn{1}{|c|}{\cellcolor[HTML]{C0C0C0}}\\
\cline{2-5}\multicolumn{1}{|c|}{\multirow{-5.25}{*}{\cellcolor[HTML]{C0C0C0}\makecell[cc]{{\footnotesize \textit{Questions sur}}\\{\footnotesize \textit{plusieurs chapitres de}}\\{\footnotesize \textit{semestres différents}}}}}&{\makecell[cc]{{\footnotesize Les métamorphoses}\\{\footnotesize du moi et histoire}\\{\footnotesize et violence}}}&\multicolumn{1}{c|}{\arabic{sujetMMHV}}&\multicolumn{1}{c}{}&\multicolumn{1}{c}{\makecell[cc]{}}&\multicolumn{1}{|c|}{\multirow{-5.5}{*}{\cellcolor[HTML]{C0C0C0}\makecell[cc]{{\footnotesize \textit{Questions qui}}\\{\footnotesize \textit{font appel à des}}\\{\footnotesize \textit{connaissances de la}}\\{\footnotesize \textit{classe de Première}}}}}\\
\hline
\end{tabularx}
\caption*{\textit{Répartition des questions selon les différents axes du programme}}
\end{table}\setlength{\itemsep}{4pt plus 2pt minus 1pt}

\section{Statistiques selon les types de sujets}

On trouvera dans les tableaux suivants le détail de la répartition des questions selon qu'elles sont des questions d'interprétation ou des essais, et selon le chapitre (ou les chapitres) auquel (ou auxquels) elles se rapportent.

Le premier tableau permet d'avoir un aperçu général des questions posées en littérature ; le deuxième permet d'avoir un aperçu général des questions posées en philosophie. On a détaillé dans ces tableaux le nombres de questions qui font appel au chapitre (à la ligne «~portant~»{}), ainsi que celle du nombre de questions pour chaque semestre (à la fin de la ligne) dans la catégorie «~Total~»{} . On a également indiqué la répartition de chacune de ces questions en-dessous de chacun des chapitres (aux lignes intitulées «~\textit{excl.}~»{}), selon qu'ils portent sur un seul chapitre du programme ou font appel à plusieurs chapitres du même semestre. On a, enfin, indiqué le détail de chacune des questions qui font appel à des sujets appartenant à deux semestres différents aux lignes intitulées «~\textit{mixtes}~»{}. La catégorie «~'Aut.'»{} désigne les questions qui font appel à des chapitres du programme de Première. Faute de place, on ne pouvait le détailler sans rendre complètement illisible le tableau ; et cela paraissait peu pertinent eu égard au faible nombre de questions dont c'est le cas puisqu'à ce jour (\today{}), seules \arabic{sujetpremterm}~questions font appel à des chapitres de Première. Certainement convient-il de noter que, même dans ce cas, la connaissance de ces chapitres est davantage utile qu'elle n'est purement nécessaire (à l'image par exemple de la connaissance utile du chapitre sur «~\HA~»{} pour traiter du sujet proposé le jour 1 en Amérique du Sud en 2022 qui propose de réfléchir à la vision de la nature qui mène à une catastrophe écologique selon Aurélien Barrau, dont le texte porte, notamment, sur la question des droits des animaux), ou encore du sujet proposé le jour 2 en Asie en 2024 pour laquelle les connaissances issues du chapitre sur «~\PdP~»{} sont utiles pour répondre à la question d'interprétation littéraire sur le texte de Rousseau (« Dans cet extrait, pourquoi le discours est-il plus apte à nous émouvoir que les images~?»{}) sans que cela soit pour autant une stricte nécessité.

Il convient de ne pas s'étonner du fait que la somme des questions portants sur un axe du programme ne corresponde pas à la somme effective de chacun des chapitres. Par exemple, s'il y a \arabic{sujetintlitteraireETA}~questions d'interprétation littéraire sur «~\ETA{}~»{}, \arabic{sujetintlitteraireEXS}~questions d'interprétation littéraire sur «~\EXS{}~»{}, \arabic{sujetintlitteraireMM}~questions d'interprétation littéraire sur «~\MM{}~»{} et \arabic{sujetintlitteraireOdS}~questions d'interprétation littéraire qui sont des «~\OdS{}~»{}, la somme totale de ces sujets n'est que de \arabic{sujetRSintlitteraire}~questions d'interprétation littéraire sur «~\RS{}»{}. Cela est dû au fait que certains sujets portant sur deux axes du même semestre sont alors comptés deux fois (par exemple les sujets sur «\EXSMM{}»{} sont comptés à la fois dans \EXS{} et dans \MM{}). La somme des questions d'interprétation littéraire sur «~\RS{}»{} est donc bien correcte. De la même manière, la somme de l'ensemble des questions littéraire (interprétation et réflexion, sur chacun des semestre) dépasse le nombre total de sujets~: cela est dû au fait que là aussi, les questions qui portent sur des chapitres appartenant à des semestres différents (par exemple les sujets sur \EXSHV{}) sont comptés deux fois dans les sommes. Si on retranche du total le nombre total de ces questions, on obtient effectivement le nombre total de questions. Chacun des tableaux suivants étant déjà très chargé en données, on a préféré omettre cette information, pour rendre l'ensemble plus lisible, d'autant plus que le nombre total de sujets recensés est déjà indiqué au début de ce chapitre de statistiques.

L'ensemble permet de se faire une idée des thèmes privilégiés selon les questions, mais il convient à ce jour (\today) de rester encore prudent~: les chapitres «~\ETA{}~»{} et «~\CC{}»{} ne pouvaient faire l'objet d'une évaluation avant la session 2024 du baccalauréat, ce qui explique l'écart en nombre de sujets et questions entre ces chapitres et les autres.

\clearpage

%\printinunitsof{cm}\prntlen{\linewidth}

\renewcommand{\arraystretch}{1}
\setlength\extrarowheight{1pt}
\begin{landscape}\label{statistiqueslittgen}
\begin{table}[h]\hypertarget{statlitt}{~}\stepcounter{subsection}\centering\phantomsection\addcontentsline{toc}{subsection}{Statistiques générales des sujets littéraires}
\begin{tabularx}{\linewidth+0.0565cm}{M{2mm}XXXM{4em}M{4em}M{2em}XXXXXM{4em}M{4em}M{2em}X}
\hline
\multicolumn{16}{c}{\makecell[cc]{{\large \textbf{Sujets littéraires}}}}\\
\hline
%\multicolumn{15}{c}{\textit{La recherche de soi}}\\
\parbox[t]{2mm}{\multirow{14}{*}{\rotatebox[origin=c]{90}{\textit{La Recherche de soi}}}}&&{Sujets}&\multicolumn{4}{c}{\footnotesize Éducation, Transmission et Émancipation}&\multicolumn{4}{c}{\footnotesize Les expressions de la sensibilité}&\multicolumn{4}{c}{\footnotesize Les métamorphoses du moi}&\multicolumn{1}{c}{\scriptsize Originaux}\\%
%%%%Début des données%%%
&\multirow{6}{*}{\makecell[cc]{\textit{Inter.}\\\textit{Littér.}}}&\textit{portant}&\multicolumn{4}{c}{\arabic{sujetallintlitteraireETA}}&\multicolumn{4}{c}{\arabic{sujetallintlitteraireEXS}}&\multicolumn{4}{c}{\arabic{sujetallintlitteraireMM}}&\arabic{sujetintlitteraireOdS}\\
\cmidrule(rl){4-7}\cmidrule(rl){8-11}\cmidrule(rl){12-15}
&&&%
ETA	&	ETAMM	&	EXSTA	&	\textit{Aut.}%
&	EXS	&	EXSTA	&	EXSMM	&	\textit{Aut.}%
&	MM	&	ETAMM	&	EXSMM	&	\textit{Aut.}%
&Total\\
&&\textit{excl.}&
\statnat{\inter}{\litte}{ETA}	& \statnat{\inter}{\litte}{ETAMM}	&	\statnat{\inter}{\litte}{EXSTA}	&	\arabic{sujetintlitterairepremtermEXS}%
&	\statnat{\inter}{\litte}{EXS}	&	\statnat{\inter}{\litte}{EXSTA}	&	\statnat{\inter}{\litte}{EXSMM}	& {\arabic{sujetintlitterairepremtermEXS}}%
&	\statnat{\inter}{\litte}{MM}	& \statnat{\inter}{\litte}{ETAMM}	&	\statnat{\inter}{\litte}{EXSMM}	&	\arabic{sujetintlitterairepremtermMM}
&\arabic{sujetRSintlitteraire}\\ \cmidrule(rl){4-7}\cmidrule(rl){8-11}\cmidrule(rl){12-15}
&&	& ETACC	&	ETAHV	&	ETAHL	&		&	EXSCC	&	EXSHV	&	EXSHL &		&	MMCC	&	MMHV	&	MMHL	&	&	\\
&&{\textit{\small mixtes}}&\statnat{\inter}{\litte}{ETACC}&\statnat{\inter}{\litte}{ETAHV}&\statnat{\inter}{\litte}{ETAHL}& &\statnat{\inter}{\litte}{EXSCC}&\statnat{\inter}{\litte}{EXSHV}&\statnat{\inter}{\litte}{EXSHL}& &\statnat{\inter}{\litte}{MMCC}&\statnat{\inter}{\litte}{MMHV}&\statnat{\inter}{\litte}{MMHL}&&\\
\cmidrule(rl){2-16}
&&&\multicolumn{4}{c}{\footnotesize Éducation, Transmission et Émancipation}&\multicolumn{4}{c}{\footnotesize Les expressions de la sensibilité}&\multicolumn{4}{c}{\footnotesize Les métamorphoses du moi}&\multicolumn{1}{c}{\scriptsize Originaux}\\
&\multirow{6}{*}{\makecell[cc]{\textit{Réfl.}\\\textit{Littér.}}}&\textit{portant}&\multicolumn{4}{c}{\arabic{sujetallreflitteraireETA}}&\multicolumn{4}{c}{\arabic{sujetallreflitteraireEXS}}&\multicolumn{4}{c}{\arabic{sujetallreflitteraireMM}}&\arabic{sujetreflitteraireOdS}\\
\cmidrule(rl){4-7}\cmidrule(rl){8-11}\cmidrule(rl){12-15}
&&&%
ETA	&	ETAMM	&	EXSTA	&	\textit{Aut.}%
&	EXS	&	EXSTA	&	EXSMM	&	\textit{Aut.}%
&	MM	&	ETAMM	&	EXSMM	&	\textit{Aut.}%
&Total\\
&&\textit{excl.}&
\statnat{\refle}{\litte}{ETA}	& \statnat{\refle}{\litte}{ETAMM}	&	\statnat{\refle}{\litte}{EXSTA}	&	\arabic{sujetintlitterairepremtermEXS}%
&	\statnat{\refle}{\litte}{EXS}	&	\statnat{\refle}{\litte}{EXSTA}	&	\statnat{\refle}{\litte}{EXSMM}	& {\arabic{sujetintlitterairepremtermEXS}}%
&	\statnat{\refle}{\litte}{MM}	& \statnat{\refle}{\litte}{ETAMM}	&	\statnat{\refle}{\litte}{EXSMM}	&	\arabic{sujetintlitterairepremtermMM}
&\arabic{sujetRSintlitteraire}\\ \cmidrule(rl){4-7}\cmidrule(rl){8-11}\cmidrule(rl){12-15}
&&	& ETACC	&	ETAHV	&	ETAHL	&		&	EXSCC	&	EXSHV	&	EXSHL &		&	MMCC	&	MMHV	&	MMHL	&	&	\\
&&{\textit{\small mixtes}}&\statnat{\refle}{\litte}{ETACC}&\statnat{\refle}{\litte}{ETAHV}&\statnat{\refle}{\litte}{ETAHL}& &\statnat{\refle}{\litte}{EXSCC}&\statnat{\refle}{\litte}{EXSHV}&\statnat{\refle}{\litte}{EXSHL}& &\statnat{\refle}{\litte}{MMCC}&\statnat{\refle}{\litte}{MMHV}&\statnat{\refle}{\litte}{MMHL}&&\\\hline
%%%%%%%%%%%%%%%%%%%%%%%%%Humanité en question%%%%%%%%%%%
%\multicolumn{15}{c}{\textit{L'Humanité en question}}\\
%\cline{3-15}
\parbox[t]{2mm}{\multirow{14}{*}{\rotatebox[origin=c]{90}{\textit{L'Humanité en question}}}}&&&\multicolumn{4}{c}{\footnotesize Création, continuités et ruptures}&\multicolumn{4}{c}{\footnotesize Histoire et violence}&\multicolumn{4}{c}{\footnotesize L'humain et ses limites}&\multicolumn{1}{c}{\scriptsize Originaux}\\%
%%%%Début des données%%%
&\multirow{6}{*}{\makecell[cc]{\textit{Inter.}\\\textit{Littér.}}}&\textit{portant}&\multicolumn{4}{c}{\arabic{sujetallintlitteraireCC}}&\multicolumn{4}{c}{\arabic{sujetallintlitteraireHV}}&\multicolumn{4}{c}{\arabic{sujetallintlitteraireHL}}&\arabic{sujetintlitteraireOdH}\\
\cmidrule(rl){4-7}\cmidrule(rl){8-11}\cmidrule(rl){12-15}
&&&%
CC	&	HVCC	&	HLCC	&	\textit{Aut.}%
&	HV	&	HVCC	&	HVL	&	\textit{Aut.}%
&	HL	&	HVL	&	HVCC	&	\textit{Aut.}%
&Total\\
&&\textit{excl.}&
\statnat{\inter}{\litte}{CC}	&	\statnat{\inter}{\litte}{HVCC}	&	\statnat{\inter}{\litte}{HLCC}	&	\arabic{sujetintlitterairepremtermCC}%
&	\statnat{\inter}{\litte}{HV}	&	\statnat{\inter}{\litte}{HVCC}	&	\statnat{\inter}{\litte}{HVL}	&	\arabic{sujetintlitterairepremtermHV}%
&	\statnat{\inter}{\litte}{HL}	&	\statnat{\inter}{\litte}{HVL}	&	\statnat{\inter}{\litte}{HLCC}	&	\arabic{sujetintlitterairepremtermHL}
&\arabic{sujetHQintlitteraire}\\ \cmidrule(rl){4-7}\cmidrule(rl){8-11}\cmidrule(rl){12-15}
&&{\multirow{2}{*}{\makecell[cc]{{\textit{{\small mixtes}}}}}}& ETACC	&	EXSCC	&	MMCC	&		&	ETAHV	&	EXSHV	&	MMHV &		&	ETAHL	&	EXSHL	&	MMHL	&	&	\\
& & & \statnat{\inter}{\litte}{ETACC}&\statnat{\inter}{\litte}{EXSCC}&\statnat{\inter}{\litte}{MMCC}& &\statnat{\inter}{\litte}{ETAHV}&\statnat{\inter}{\litte}{EXSHV}&\statnat{\inter}{\litte}{MMHV}& &\statnat{\inter}{\litte}{ETAHL}&\statnat{\inter}{\litte}{EXSHL}&\statnat{\inter}{\litte}{MMHL}&&\\
\cmidrule(rl){2-16}
&&&\multicolumn{4}{c}{\footnotesize Création, continuités et ruptures}&\multicolumn{4}{c}{\footnotesize Histoire et violenec}&\multicolumn{4}{c}{\footnotesize L'humain et ses limites}&\multicolumn{1}{c}{\scriptsize Originaux}\\%
&\multirow{6}{*}{\makecell[cc]{\textit{Réfl.}\\\textit{Littér.}}}&\textit{portant}&\multicolumn{4}{c}{\arabic{sujetallreflitteraireCC}}&\multicolumn{4}{c}{\arabic{sujetallreflitteraireHV}}&\multicolumn{4}{c}{\arabic{sujetallreflitteraireHL}}&\arabic{sujetreflitteraireOdH}\\
\cmidrule(rl){4-7}\cmidrule(rl){8-11}\cmidrule(rl){12-15}
&&&%
CC	&	HVCC	&	HLCC	&	\textit{Aut.}%
&	HV	&	HVCC	&	HVL	&	\textit{Aut.}%
&	HL	&	HVL	&	HVCC	&	\textit{Aut.}%
&Total\\
&&\textit{excl.}&
\statnat{\refle}{\litte}{CC}	&	\statnat{\refle}{\litte}{HVCC}	&	\statnat{\refle}{\litte}{HLCC}	&	\arabic{sujetintlitterairepremtermCC}%
&	\statnat{\refle}{\litte}{HV}	&	\statnat{\refle}{\litte}{HVCC}	&	\statnat{\refle}{\litte}{HVL}	&	\arabic{sujetintlitterairepremtermHV}%
&	\statnat{\refle}{\litte}{HL}	&	\statnat{\refle}{\litte}{HVL}	&	\statnat{\refle}{\litte}{HLCC}	&	\arabic{sujetintlitterairepremtermHL}
&\arabic{sujetHQintlitteraire}\\ \cmidrule(rl){4-7}\cmidrule(rl){8-11}\cmidrule(rl){12-15}
&&{\multirow{2}{*}{\makecell[cc]{{\textit{{\small mixtes}}}}}}& ETACC	&	EXSCC	&	MMCC	&		&	ETAHV	&	EXSHV	&	MMHV &		&	ETAHL	&	EXSHL	&	MMHL	&	&	\\
& & & \statnat{\refle}{\litte}{ETACC}&\statnat{\refle}{\litte}{EXSCC}&\statnat{\refle}{\litte}{MMCC}& &\statnat{\refle}{\litte}{ETAHV}&\statnat{\refle}{\litte}{EXSHV}&\statnat{\refle}{\litte}{MMHV}& &\statnat{\refle}{\litte}{ETAHL}&\statnat{\refle}{\litte}{EXSHL}&\statnat{\refle}{\litte}{MMHL}&&\\\hline
\multicolumn{16}{c}{\makecell[cc]{{\scriptsize \underline{Liste des abréviations :} ETA~: \ETA{}, EXS~: \EXS{}, MM~: \MM{}, CC~: \CC,}\\%
{\scriptsize HV~: \HV{}, HL~: \HL{}, ETAMM~: \ETAMM, EXSTA~: \EXS}\\{\scriptsize et \ETA{}, HVL~: \HVL, EXSMM~: \EXSMM, ETACC~:}\\{\scriptsize \ETACC, ETAHV~: \ETAHV{}, ETAHL~: Éducation, transmission et}\\{\scriptsize émancipation et \HL, EXSCC~: \EXSCC, EXSHV~: \EXSHV, EXSHL~:}\\{\scriptsize  \EXSHL, MMCC~: \MMCC, MMHV~: \MMHV}\\{\scriptsize MMHL~: \MMHL. Dans la catégorie '\textit{Aut}', on trouve les sujets en lien avec la classe de Première.}}}\\
\end{tabularx}
\end{table}
\end{landscape}

\clearpage
\begin{landscape}
\begin{table}[h]\stepcounter{subsection}\centering\phantomsection\addcontentsline{toc}{subsection}{Statistiques générales des sujets philosophiques}
\begin{tabularx}{\linewidth+0.0565cm}{M{2mm}XXXM{4em}M{4em}M{2em}XXXXXM{4em}M{4em}M{2em}X}
\hline
\multicolumn{16}{c}{\makecell[cc]{{\large \textbf{Sujets philosophiques}}}}\\
\hline
%\multicolumn{15}{c}{\textit{La recherche de soi}}\\
\parbox[t]{2mm}{\multirow{14}{*}{\rotatebox[origin=c]{90}{\textit{La Recherche de soi}}}}&&{Sujets}&\multicolumn{4}{c}{\footnotesize Éducation, Transmission et Émancipation}&\multicolumn{4}{c}{\footnotesize Les expressions de la sensibilité}&\multicolumn{4}{c}{\footnotesize Les métamorphoses du moi}&\multicolumn{1}{c}{\scriptsize Originaux}\\%
%%%%Début des données%%%
&\multirow{6}{*}{\makecell[cc]{\textit{Inter.}\\\textit{Phil.}}}&\textit{portant}&\multicolumn{4}{c}{\arabic{sujetallintphilosophiqueETA}}&\multicolumn{4}{c}{\arabic{sujetallintphilosophiqueEXS}}&\multicolumn{4}{c}{\arabic{sujetallintphilosophiqueMM}}&\arabic{sujetintphilosophiqueOdS}\\
\cmidrule(rl){4-7}\cmidrule(rl){8-11}\cmidrule(rl){12-15}
&&&%
ETA	&	ETAMM	&	EXSTA	&	\textit{Aut.}%
&	EXS	&	EXSTA	&	EXSMM	&	\textit{Aut.}%
&	MM	&	ETAMM	&	EXSMM	&	\textit{Aut.}%
&Total\\
&&\textit{excl.}&
\statnat{\inter}{\phil}{ETA}	& \statnat{\inter}{\phil}{ETAMM}	&	\statnat{\inter}{\phil}{EXSTA}	&	\arabic{sujetintphilosophiquepremtermEXS}%
&	\statnat{\inter}{\phil}{EXS}	&	\statnat{\inter}{\phil}{EXSTA}	&	\statnat{\inter}{\phil}{EXSMM}	& {\arabic{sujetintphilosophiquepremtermEXS}}%
&	\statnat{\inter}{\phil}{MM}	& \statnat{\inter}{\phil}{ETAMM}	&	\statnat{\inter}{\phil}{EXSMM}	&	\arabic{sujetintphilosophiquepremtermMM}
&\arabic{sujetRSintphilosophique}\\ \cmidrule(rl){4-7}\cmidrule(rl){8-11}\cmidrule(rl){12-15}
&&	& ETACC	&	ETAHV	&	ETAHL	&		&	EXSCC	&	EXSHV	&	EXSHL &		&	MMCC	&	MMHV	&	MMHL	&	&	\\
&&{\textit{\small mixtes}}&\statnat{\inter}{\phil}{ETACC}&\statnat{\inter}{\phil}{ETAHV}&\statnat{\inter}{\phil}{ETAHL}& &\statnat{\inter}{\phil}{EXSCC}&\statnat{\inter}{\phil}{EXSHV}&\statnat{\inter}{\phil}{EXSHL}& &\statnat{\inter}{\phil}{MMCC}&\statnat{\inter}{\phil}{MMHV}&\statnat{\inter}{\phil}{MMHL}&&\\
\cmidrule(rl){2-16}
&&&\multicolumn{4}{c}{\footnotesize Éducation, Transmission et Émancipation}&\multicolumn{4}{c}{\footnotesize Les expressions de la sensibilité}&\multicolumn{4}{c}{\footnotesize Les métamorphoses du moi}&\multicolumn{1}{c}{\scriptsize Originaux}\\
&\multirow{6}{*}{\makecell[cc]{\textit{Réfl.}\\\textit{Phil.}}}&\textit{portant}&\multicolumn{4}{c}{\arabic{sujetallrefphilosophiqueETA}}&\multicolumn{4}{c}{\arabic{sujetallrefphilosophiqueEXS}}&\multicolumn{4}{c}{\arabic{sujetallrefphilosophiqueMM}}&\arabic{sujetrefphilosophiqueOdS}\\
\cmidrule(rl){4-7}\cmidrule(rl){8-11}\cmidrule(rl){12-15}
&&&%
ETA	&	ETAMM	&	EXSTA	&	\textit{Aut.}%
&	EXS	&	EXSTA	&	EXSMM	&	\textit{Aut.}%
&	MM	&	ETAMM	&	EXSMM	&	\textit{Aut.}%
&Total\\
&&\textit{excl.}&
\statnat{\refle}{\phil}{ETA}	& \statnat{\refle}{\phil}{ETAMM}	&	\statnat{\refle}{\phil}{EXSTA}	&	\arabic{sujetintphilosophiquepremtermEXS}%
&	\statnat{\refle}{\phil}{EXS}	&	\statnat{\refle}{\phil}{EXSTA}	&	\statnat{\refle}{\phil}{EXSMM}	& {\arabic{sujetintphilosophiquepremtermEXS}}%
&	\statnat{\refle}{\phil}{MM}	& \statnat{\refle}{\phil}{ETAMM}	&	\statnat{\refle}{\phil}{EXSMM}	&	\arabic{sujetintphilosophiquepremtermMM}
&\arabic{sujetRSintphilosophique}\\ \cmidrule(rl){4-7}\cmidrule(rl){8-11}\cmidrule(rl){12-15}
&&	& ETACC	&	ETAHV	&	ETAHL	&		&	EXSCC	&	EXSHV	&	EXSHL &		&	MMCC	&	MMHV	&	MMHL	&	&	\\
&&{\textit{\small mixtes}}&\statnat{\refle}{\phil}{ETACC}&\statnat{\refle}{\phil}{ETAHV}&\statnat{\refle}{\phil}{ETAHL}& &\statnat{\refle}{\phil}{EXSCC}&\statnat{\refle}{\phil}{EXSHV}&\statnat{\refle}{\phil}{EXSHL}& &\statnat{\refle}{\phil}{MMCC}&\statnat{\refle}{\phil}{MMHV}&\statnat{\refle}{\phil}{MMHL}&&\\\hline
%%%%%%%%%%%%%%%%%%%%%%%%%Humanité en question%%%%%%%%%%%
%\multicolumn{15}{c}{\textit{L'Humanité en question}}\\
%\cline{3-15}
\parbox[t]{2mm}{\multirow{14}{*}{\rotatebox[origin=c]{90}{\textit{L'Humanité en question}}}}&&&\multicolumn{4}{c}{\footnotesize Création, continuités et ruptures}&\multicolumn{4}{c}{\footnotesize Histoire et violence}&\multicolumn{4}{c}{\footnotesize L'humain et ses limites}&\multicolumn{1}{c}{\scriptsize Originaux}\\%
%%%%Début des données%%%
&\multirow{6}{*}{\makecell[cc]{\textit{Inter.}\\\textit{Phil.}}}&\textit{portant}&\multicolumn{4}{c}{\arabic{sujetallintphilosophiqueCC}}&\multicolumn{4}{c}{\arabic{sujetallintphilosophiqueHV}}&\multicolumn{4}{c}{\arabic{sujetallintphilosophiqueHL}}&\arabic{sujetintphilosophiqueOdH}\\
\cmidrule(rl){4-7}\cmidrule(rl){8-11}\cmidrule(rl){12-15}
&&&%
CC	&	HVCC	&	HLCC	&	\textit{Aut.}%
&	HV	&	HVCC	&	HVL	&	\textit{Aut.}%
&	HL	&	HVL	&	HVCC	&	\textit{Aut.}%
&Total\\
&&\textit{excl.}&
\statnat{\inter}{\phil}{CC}	&	\statnat{\inter}{\phil}{HVCC}	&	\statnat{\inter}{\phil}{HLCC}	&	\arabic{sujetintphilosophiquepremtermCC}%
&	\statnat{\inter}{\phil}{HV}	&	\statnat{\inter}{\phil}{HVCC}	&	\statnat{\inter}{\phil}{HVL}	&	\arabic{sujetintphilosophiquepremtermHV}%
&	\statnat{\inter}{\phil}{HL}	&	\statnat{\inter}{\phil}{HVL}	&	\statnat{\inter}{\phil}{HLCC}	&	\arabic{sujetintphilosophiquepremtermHL}
&\arabic{sujetHQintphilosophique}\\ \cmidrule(rl){4-7}\cmidrule(rl){8-11}\cmidrule(rl){12-15}
&&{\multirow{2}{*}{\makecell[cc]{{\textit{{\small mixtes}}}}}}& ETACC	&	EXSCC	&	MMCC	&		&	ETAHV	&	EXSHV	&	MMHV &		&	ETAHL	&	EXSHL	&	MMHL	&	&	\\
& & & \statnat{\inter}{\phil}{ETACC}&\statnat{\inter}{\phil}{EXSCC}&\statnat{\inter}{\phil}{MMCC}& &\statnat{\inter}{\phil}{ETAHV}&\statnat{\inter}{\phil}{EXSHV}&\statnat{\inter}{\phil}{MMHV}& &\statnat{\inter}{\phil}{ETAHL}&\statnat{\inter}{\phil}{EXSHL}&\statnat{\inter}{\phil}{MMHL}&&\\
\cmidrule(rl){2-16}
&&&\multicolumn{4}{c}{\footnotesize Création, continuités et ruptures}&\multicolumn{4}{c}{\footnotesize Histoire et violenec}&\multicolumn{4}{c}{\footnotesize L'humain et ses limites}&\multicolumn{1}{c}{\scriptsize Originaux}\\%
&\multirow{6}{*}{\makecell[cc]{\textit{Réfl.}\\\textit{Phil.}}}&\textit{portant}&\multicolumn{4}{c}{\arabic{sujetallrefphilosophiqueCC}}&\multicolumn{4}{c}{\arabic{sujetallrefphilosophiqueHV}}&\multicolumn{4}{c}{\arabic{sujetallrefphilosophiqueHL}}&\arabic{sujetrefphilosophiqueOdH}\\
\cmidrule(rl){4-7}\cmidrule(rl){8-11}\cmidrule(rl){12-15}
&&&%
CC	&	HVCC	&	HLCC	&	\textit{Aut.}%
&	HV	&	HVCC	&	HVL	&	\textit{Aut.}%
&	HL	&	HVL	&	HVCC	&	\textit{Aut.}%
&Total\\
&&\textit{excl.}&
\statnat{\refle}{\phil}{CC}	&	\statnat{\refle}{\phil}{HVCC}	&	\statnat{\refle}{\phil}{HLCC}	&	\arabic{sujetintphilosophiquepremtermCC}%
&	\statnat{\refle}{\phil}{HV}	&	\statnat{\refle}{\phil}{HVCC}	&	\statnat{\refle}{\phil}{HVL}	&	\arabic{sujetintphilosophiquepremtermHV}%
&	\statnat{\refle}{\phil}{HL}	&	\statnat{\refle}{\phil}{HVL}	&	\statnat{\refle}{\phil}{HLCC}	&	\arabic{sujetintphilosophiquepremtermHL}
&\arabic{sujetHQintphilosophique}\\ \cmidrule(rl){4-7}\cmidrule(rl){8-11}\cmidrule(rl){12-15}
&&{\multirow{2}{*}{\makecell[cc]{{\textit{{\small mixtes}}}}}}& ETACC	&	EXSCC	&	MMCC	&		&	ETAHV	&	EXSHV	&	MMHV &		&	ETAHL	&	EXSHL	&	MMHL	&	&	\\
& & & \statnat{\refle}{\phil}{ETACC}&\statnat{\refle}{\phil}{EXSCC}&\statnat{\refle}{\phil}{MMCC}& &\statnat{\refle}{\phil}{ETAHV}&\statnat{\refle}{\phil}{EXSHV}&\statnat{\refle}{\phil}{MMHV}& &\statnat{\refle}{\phil}{ETAHL}&\statnat{\refle}{\phil}{EXSHL}&\statnat{\refle}{\phil}{MMHL}&&\\\hline
\multicolumn{16}{c}{\makecell[cc]{{\scriptsize \underline{Liste des abréviations :} ETA~: \ETA{}, EXS~: \EXS{}, MM~: \MM{}, CC~: \CC,}\\%
{\scriptsize HV~: \HV{}, HL~: \HL{}, ETAMM~: \ETAMM, EXSTA~: \EXS}\\{\scriptsize et \ETA{}, HVL~: \HVL, EXSMM~: \EXSMM, ETACC~:}\\{\scriptsize \ETACC, ETAHV~: \ETAHV{}, ETAHL~: Éducation, transmission et}\\{\scriptsize émancipation et \HL, EXSCC~: \EXSCC, EXSHV~: \EXSHV, EXSHL~:}\\{\scriptsize  \EXSHL, MMCC~: \MMCC, MMHV~: \MMHV}\\{\scriptsize MMHL~: \MMHL. Dans la catégorie '\textit{Aut}', on trouve les sujets en lien avec la classe de Première et les questions « originales ».}}}\\
\end{tabularx}
\end{table}
\end{landscape}

\begin{landscape}
\statexam{Général}
\end{landscape}

\begin{landscape}
\graphexamen{l'ensemble des centres}
\end{landscape}


\newpage

\section{Statistiques par centres d'examen}
%\phantomsection\addcontentsline{toc}{section}{Statistiques par centres d'examen}

De manière à disposer d'une vision plus granulaire selon les centres d'examen, on a jugé utile de proposer des statistiques par centres d'examen : Asie, Amérique du Nord, Centres étrangers, Métropole, Nouvelle-Calédonie, Polynésie\footnote{On n'a pas indiqué spécifiquement les sujets pour l'Amérique du Sud, les Antilles, le Liban la Réunion en raison du faible nombre de sujets proposés~: ces centres d'examen ont rejoint, au cours des années, d'autres centres d'examen, s'en sont dissociés (Antilles, 2026) ou n'ont pas proposé régulièrement de sujets d'Humanités (particulièrement l'Amérique du Sud).}.

On a essayé de rendre le plus lisible possible ces statistique, mais au vu du nombre de possibilités laissées par les questions selon le chapitre auquel elles se rattachent, la possibilité d'avoir des sujets croisés tant au sein d'un même semestre qu'avec l'autre semestre, voire avec des sujets qui peuvent faire appel à des connaissances de la classe de Première, cela reste nécessairement un peu complexe. On a laissé, parce que cela nous paraît le plus important, dans le premier tableau, du haut, tous les sujets relatifs aux intitulés les plus «~classique~» de Terminale, soit que les questions portent sur une seule notion, soit qu'ils croisent des notions d'un même semestre. On a préservé la séparation entre le semestre portant sur la Recherche de soi et l'Humanité en question, et au sein de chacun des ses thèmes de semestre, on a distingué les sujets qui font exclusivement appel à un chapitre, et ceux qui font appel à deux chapitres du même semestre. Dans le deuxième tableau, en dessous du précédent, on trouve à gauche les questions qui font appel à des chapitres de deux semestres différents. Dans le dernier tableau, en bas de page, on a réindiqué le nombre total de sujet selon les différents chapitres, en comptant les sujets qui portent exclusivement sur un chapitre et ceux qui font appel à plusieurs chapitres. Cela permet de se faire une bonne idée, je crois, de la manière dont ont été répartis les sujets dans chaque centre d'examen. On a indiqué, pour la lisibilité, l'ensemble des abréviations utilisés en bas à droite de chaque page, et le nombre de sujets proposés au total par centre d'examens.

Dans la mesure où, par centres d'examens, le nombre total de sujet n'est pas si élevé que cela, on voit nécessairement apparaître un nombre important d'absences de sujets proposés. Avec les années, ces absences devraient probablement disparaître. Les sujets qui font appel à des connaissances de Première étant plus rares, c'est le tableau qui risque de rester le plus «~vide~», et l'on verra, à l'usage, si on le conserve. Pour l'heure, il n'y a qu'en Amérique du Sud, au Liban et une fois en Asie que de tels sujets ont été proposés.

Comme pour les autres statistiques, la présence de questions qui portent sur des chapitres croisés, soit au sein d'un même semestre, soit au sein de différents semestres, soit même au sein de différentes années rend impossible le fait de seulement faire la somme des questions indiquées dans «~Tous les axes confondus~». C'est la raison pour laquelle on a indiqué le nombre de sujets par centre d'examen en bas de page~: on se fait ainsi une meilleure idée des sujets proposés en ce qu'on peut déterminer le \textit{ratio} entre le nombre de sujets qui portent sur un axe et le nombre total de sujets.

\vfill{~}\newpage

\newpage

\begin{landscape}
\statexam[amnord]{Amérique du Nord}

\graphexamen[amnord]{l'Amérique du Nord}
\end{landscape}

\newpage

\begin{landscape}
\statexam[asi]{Asie}

\graphexamen[asi]{l'Asie}

\end{landscape}

\begin{landscape}
\statexam[cea]{Centres étrangers gr. 1}

\graphexamen[cea]{les Centres étrangers (gr.~1)}
\end{landscape}

\newpage

\begin{landscape}
\statexam[met]{Métropole}

\graphexamen[met]{la Métropole}
\end{landscape}

\newpage

\begin{landscape}
\statexam[nca]{Nouvelle-Calédonie}

\graphexamen[nca]{la Nouvelle-Calédonie}
\end{landscape}

\newpage

\begin{landscape}
\statexam[pol]{Polynésie}

\graphexamen[pol]{la Polynésie}
\end{landscape}